\documentclass[12pt,a4paper]{article}
\usepackage[margin=1in]{geometry}
\usepackage{amsmath,amssymb,amsthm}
\usepackage{hyperref}
\usepackage{enumitem}
\usepackage{booktabs}
\usepackage{array}

\newtheorem{theorem}{Theorem}[section]
\newtheorem{corollary}[theorem]{Corollary}
\newtheorem{lemma}[theorem]{Lemma}
\newtheorem{proposition}[theorem]{Proposition}
\newtheorem{definition}[theorem]{Definition}
\newtheorem*{remark}{Remark}

\title{\textbf{The Cascade Series --- Part VI}\\[6pt]
Tower Growth, Inflation, and the Pre-Big-Bang\\
from the Cascade}
\author{RTAC}
\date{April 2026}

\begin{document}
\maketitle

\begin{center}
\fbox{\parbox{0.9\textwidth}{\textbf{Status.} Speculative extension. Tier 5 throughout unless otherwise noted. This paper proposes a dynamical reading of the cascade's descent that is not required by Parts~0--V. It is included as a forward extension and a source of falsifiable predictions; it does not feed back into any Tier 1--3 result of the preceding papers.}}
\end{center}

\begin{abstract}
Parts~0--V~\cite{paper0,paper1,paper2,paper3,paper4a,paper4b,paper5} treat the cascade's descent from $d = \infty$ to $d = 4$ as a timeless logical relationship. This paper promotes the truncation height to a physical time coordinate: at time $t$, the cascade above the observer is truncated at $N(t) = N_0 + t/(\alpha\,t_{\mathrm{Pl,red}})$, with one Planck tick adding one layer at the top. The observer's dimension remains fixed at $d = 4$.

Under this reading: (1)~the pre-Big-Bang universe has four discrete phases (A: $N=4$--$6$, B: $N=7$--$18$, C: $N=19$--$216$, D: $N \geq 217$), each with cascade-determined $\rho_\Lambda$ and Hubble rate. (2)~Total e-folds during phases A--C: $N_e = 70.94$ at tick rate $\alpha = 1$, inside the observational inflation window. (3)~The Big Bang is the $N=217$ phase transition: $\Omega_{217}$ activates in one Planck tick, $\rho_\Lambda$ drops by $10^{120}$, $\Delta\rho \approx 0.3\,M_{\mathrm{Pl,red}}^4$ thermalises to radiation at $T_{\mathrm{RH}} \sim 7 \times 10^{17}$~GeV. (4)~The cascade's native perturbation source is the Gram overlap deficit $\sqrt{1 - C_{d,d+1}^2}$ between consecutive layers. (5)~Tensor modes are structurally suppressed: the metric is a state property (Part II\,=\,III~\cite{paper23}), no quantised graviton exists (Part~III~\cite{paper3}~\S12). (6)~The Standard Model's structural features assemble in a specific cascade-internal order as the tower grows.

Falsifiers: CMB-S4 $n_s$ functional form; LiteBIRD tensor detection at $r \gtrsim 10^{-3}$; deviations from the cascade-native particle-structural timeline.
\end{abstract}

\tableofcontents
\newpage

\section{Division of Labour}

Papers~0--V derive the cascade's timeless structural content. This paper overlays a dynamical reading in which the cascade's truncation height evolves with cosmic time. Part~VI does not modify any conclusion of Parts~0--V; it extends them.

\begin{center}
\begin{tabular}{lll}
\toprule
Paper & Provides & Status \\
\midrule
0~\cite{paper0} & Four distinguished dimensions, invariant $I$ & Theorem \\
I~\cite{paper1} & $\rho_\Lambda/M_{\mathrm{Pl,red}}^4$ from observer frame & Tier 2 \\
II~\cite{paper2} & QM from cascade; time = slicing direction & Tier 1--2 \\
III~\cite{paper3} & GR, $d = 4$, Lorentzian signature & Tier 1 \\
II\,=\,III~\cite{paper23} & Quantum gravity dissolution & Tier 1 \\
IVa--b~\cite{paper4a,paper4b} & SM group, masses, couplings & Tier 1--3 \\
V~\cite{paper5} & Cosmological parameters today & Tier 2--3 \\
\textbf{VI (this)} & \textbf{Tower growth, cosmic history, inflation} & \textbf{Tier 5} \\
\bottomrule
\end{tabular}
\end{center}

\section{Time as Tower Growth}

\subsection{The identification}

\begin{definition}[Tower-growth time]\label{def:tower-time}
At time $t$, the cascade above the observer at $d = 4$ is truncated at
\[
N(t) \;=\; N_0 \,+\, \frac{t}{\alpha\,t_{\mathrm{Pl,red}}},
\]
where $\alpha$ is a dimensionless tick rate, $N_0$ is an initial truncation
height, and $t_{\mathrm{Pl,red}}$ is the reduced Planck time. One Planck tick
adds one layer at the top of the tower.
\end{definition}

The observer's dimension $d = 4$ remains fixed at all times. What changes with
time is the \emph{height} of the structurally-present tower above the observer.
Physical cosmic time is the tower's growth, not a coordinate the cascade lives
in.

\subsection{Reconciliation with Paper~II}

Paper~II~\cite{paper2}, Section~7.1 identifies time with the cascade's slicing
direction. The tower-growth reading refines that identification into two coupled
operations:

\begin{itemize}[nosep]
\item \textbf{Within-moment projection.} At fixed $t$, the map from $d = N(t)$
down to $d = 4$ through the currently-existing tower. This is Paper~II's
discrete propagator $K(N(t), 4) = \prod_{j = 4}^{N(t) - 1} L(j)$.

\item \textbf{Between-moments evolution.} $N(t) \to N(t) + 1$: a new layer is
added at the top of the tower. This operation has no analogue in Paper~II.
\end{itemize}

\begin{proposition}[Physical time evolution]\label{prop:time-evolution}
Physical evolution of the cascade state from time $t$ to $t + \alpha\,
t_{\mathrm{Pl,red}}$ is the composition
\[
|\Psi(t + \alpha\,t_{\mathrm{Pl,red}})\rangle
\;=\; K(N(t) + 1, 4) \;\circ\; A_{N(t)+1} \;\circ\; K(N(t), 4)^{-1} \;
|\Psi(t)\rangle,
\]
where $A_{N(t)+1}$ is the tower-growth operator that adds layer $N(t) + 1$ to
the tower. Paper~II's discrete propagator $K$ alone is a static within-moment
projection; the tower-growth step $A$ is the genuine dynamical ingredient.
\end{proposition}

\begin{remark}
This resolves a latent tension in Paper~II. The discrete propagator $K(D, d')$
is presented there as ``the cascade propagator'', but it acts between fixed
cascade dimensions rather than between moments of time. Under Part~VI's
identification, $K$ is the static projection at each moment, and the evolution
between moments is the addition of one layer at the top followed by
re-projection. Both operations are cascade-internal; Paper~II captures the
first, Part~VI the second.
\end{remark}

\subsection{Candidate tick rates}

Two cascade-native values for $\alpha$ are available:

\begin{center}
\renewcommand{\arraystretch}{1.25}
\begin{tabular}{llcc}
\toprule
Reading & Tick = & Phase~C duration & $N_e^{\mathrm{total}}$ \\
\midrule
(i) & $t_{\mathrm{Pl,red}}$ & $198\,t_{\mathrm{Pl,red}}$ & $70.94$ \\
(ii) & $(3\pi/8)\,t_{\mathrm{Pl,red}}$ (cascade lapse $N(4)$) &
      $233\,t_{\mathrm{Pl,red}}$ & $81.6$ \\
\bottomrule
\end{tabular}
\end{center}

Reading~(i) uses the reduced Planck time directly. Reading~(ii) uses the
cascade's own lapse at the observer's dimension, $N(4) = 3\pi/8$ (Paper~II~\S7.2,
from the slicing recurrence at $d = 4$). Both values are derived from cascade
structure; the criterion for selecting between them is open (Section~\ref{sec:open}).
Numerical results below use reading~(i), which matches the simulator
implementation of Issue~\#65.

\subsection{Initial truncation height}\label{subsec:N0}

The choice of the initial truncation height $N_0$ is not forced by Part~VI;
it depends on which identification of the observer is taken as primary. The
two candidates rest on different readings of the same $S^3$:

\begin{itemize}[leftmargin=*,nosep]
\item \textbf{Paper~II reading:} The observer at $d = 4$ has state space
$S^{d-1} = S^3 = \partial B^4$ by boundary dominance
(Paper~II~\cite{paper2}~\S2.3). The observer's 3-sphere is the boundary of
the 4-ball; only $d = 4$ is structurally required. This gives $N_0 = 4$.
\item \textbf{Cover-sheet reading:} The observer lives on the $S^3$ horizon
of a 5D black hole (Cover sheet, ``The Thought Experiment''). That $S^3$ is
the horizon of a 5-dimensional black hole and requires $d = 5$ to be
structurally present. This gives $N_0 = 5$.
\end{itemize}

The two readings agree on what the observer \emph{is} (a 3-sphere) but differ
on what minimum host structure is required for the 3-sphere to exist. The
$N_0 = 4$ reading treats the observer's $S^3$ as a freestanding boundary of
its own 4-ball; the $N_0 = 5$ reading treats it as a horizon embedded in a
5D ambient geometry. Paper~I~\cite{paper1}~\S3.2 uses the cover-sheet reading
to justify the locked-time 3D projection factor, favouring $N_0 = 5$ for
consistency; Paper~II's boundary-dominance derivation is agnostic.

The simulator of Issue~\#65 uses $N_0 = 4$. Under $N_0 = 5$, Phase~A
shortens from 3 ticks to 2, reducing total e-folds by $\Delta N_e \approx
-0.46$ (one Phase~A tick of $H = 0.4607$). Selection between the two is
open; Section~\ref{sec:open}, Open Question~2.

\section{Phase Structure of the Pre-Big-Bang Universe}
\label{sec:phases}

\subsection{Factor activation}

At each distinguished cascade layer, an additional sphere-area factor enters the
cosmological constant ratio $\rho_\Lambda / M_{\mathrm{Pl,red}}^4$. Following
Paper~I~\cite{paper1} Theorem~3.1 and Paper~V~\cite{paper5} Section~2, the
cumulative factor at truncation height $N$ is the product of all distinguished
contributions at or below $N$:

\begin{center}
\renewcommand{\arraystretch}{1.25}
\begin{tabular}{ccl}
\toprule
Layer & Factor & Source \\
\midrule
$N \geq 4$ & $2/\pi$ & Observer locked-time 3D projection (Paper~I \S3.2) \\
$N \geq 7$ & $(\Omega_5/\Omega_7)^2 = 9/\pi^2$ & Host frame $d_V \to d_0$ (Paper~I \S3.1) \\
$N \geq 19$ & $\Omega_{19} \approx 0.51614$ & First threshold $d_1$ (Paper~0 Thm~6.7) \\
$N \geq 217$ & $\Omega_{217} \approx 2.335 \times 10^{-120}$ & Second threshold $d_2$ \\
\bottomrule
\end{tabular}
\end{center}

Each factor switches on the first tick at which its associated layer enters the
truncated tower. Before activation the factor is absent; after activation it is
present in all subsequent ticks. The activation is discrete (single-tick),
forcing the phase structure below.

\subsection{Phase table}

\begin{center}
\renewcommand{\arraystretch}{1.3}
\begin{tabular}{clcccc}
\toprule
Phase & $N$ range & $\rho_\Lambda / M_{\mathrm{Pl,red}}^4$ & $H / M_{\mathrm{Pl,red}}$ & Ticks & E-folds \\
\midrule
A & $4$--$6$ & $2/\pi = 0.6366$ & $0.4607$ & $3$ & $1.38$ \\
B & $7$--$18$ & $(2/\pi)(9/\pi^2) = 0.5805$ & $0.4420$ & $12$ & $5.30$ \\
C & $19$--$216$ & Phase~B $\times\,\Omega_{19} = 0.2996$ & $0.3168$ & $198$ & $62.73$ \\
D & $\geq 217$ & Phase~C $\times\,\Omega_{217} = 7.00 \times 10^{-121}$ & $\sim 10^{-61}$ & $\sim 10^{60}$ & $1.52$ + expansion \\
\bottomrule
\end{tabular}
\end{center}

Hubble rates follow from the standard Friedmann equation
$H^2 = \rho_\Lambda / (3 \Omega_\Lambda M_{\mathrm{Pl,red}}^2)$ applied with the
cascade's phase-specific $\rho_\Lambda$ under $\Omega_\Lambda \approx 1$ (vacuum
domination during inflation). Within-phase evolution is pure de~Sitter
($H$ approximately constant); transitions between phases are single-tick
discontinuities where the next factor activates.

\subsection{E-fold count}

\begin{proposition}[Pre-Big-Bang e-folds]\label{prop:efolds}
Under tick-rate reading~(i), the total number of e-folds accumulated from
$N = N_0$ to $N = 217$ is $N_e = 70.94$, inside the observational inflation
window $[50, 70]$ (with mild extension at the upper end for non-slow-roll
regimes).
\end{proposition}

\begin{proof}
Each tick contributes $H(\mathrm{phase}) \cdot t_{\mathrm{Pl,red}}$ e-folds.
Summing across the four pre-Big-Bang phases:
\begin{align*}
N_e^{A} &= 3 \times 0.4607 = 1.382, \\
N_e^{B} &= 12 \times 0.4420 = 5.304, \\
N_e^{C} &= 198 \times 0.3168 = 62.732, \\
N_e^{D, \mathrm{early}} &= 1.523 \text{ (first 30 ticks of Phase~D before radiation dominates).}
\end{align*}
Total pre-radiation e-folds: $N_e = 70.94$. Verification by the numerical
simulator of Issue~\#65 (\texttt{tools/generators/tower\_growth\_simulator.py}) confirms
this value to four decimal places.
\end{proof}

\begin{remark}[Phase~C dominance]
The 63 e-folds contributed by Phase~C come from $198$ Planck ticks at
$H/M_{\mathrm{Pl,red}} \approx 0.317$. The phase duration is structurally
forced: $216 - 19 + 1 = 198$ is the number of cascade layers between the two
thresholds of Paper~0. The factor of approximately $60$ e-folds required to
solve the horizon and flatness problems is not chosen to match observation; it
is fixed by Paper~0's bilinear-invariant structure $217 = 7 \times 31$ and the
Phase~C Hubble rate determined by the cascade's cosmological-constant product
up to $d_1$.
\end{remark}

\begin{remark}[Why the window is hit, structurally]
Two independent cascade numbers produce the 70 e-folds: the phase-C Hubble
rate $H/M_{\mathrm{Pl,red}}$ is set by $\sqrt{(2/\pi)(9/\pi^2)\Omega_{19}/3}$
(Paper~0 and Paper~I quantities), and the phase-C duration is set by
$217 - 19 = 198$ ticks (Paper~0 threshold spacing). Their product lands inside
the observational window without fitting. That this ``lands'' is either a
genuine structural prediction or a suggestive coincidence; Section~\ref{sec:open}
lists the tightening required to distinguish.
\end{remark}

\section{The $N = 217$ Big Bang}
\label{sec:bigbang}

\subsection{The discrete phase transition}

At $N = 217$, the activation of $\Omega_{217} \approx 2.33 \times 10^{-120}$
multiplies the cosmological constant ratio by that factor in a single Planck
tick. Before the tick (Phase~C): $\rho_\Lambda / M_{\mathrm{Pl,red}}^4 \approx
0.30$. After the tick (Phase~D): $\rho_\Lambda / M_{\mathrm{Pl,red}}^4 \approx
7 \times 10^{-121}$. The drop is $120$ orders of magnitude in one
$t_{\mathrm{Pl,red}}$.

\begin{theorem}[Discrete reheating]\label{thm:reheating}
At the $N = 217$ tick, the cosmological-constant energy density
$\rho_\Lambda^{(C)} \approx 0.30\,M_{\mathrm{Pl,red}}^4$ is replaced by
$\rho_\Lambda^{(D)} \approx 7 \times 10^{-121}\,M_{\mathrm{Pl,red}}^4$. The
released energy
\[
\Delta\rho \;=\; \rho_\Lambda^{(C)} - \rho_\Lambda^{(D)}
\;\approx\; 0.30\,M_{\mathrm{Pl,red}}^4
\]
thermalises into the cascade's $w = 1/3$ metric content (Paper~III~\cite{paper3}
Theorem~14.6), producing a radiation-dominated early universe with no matter
content.
\end{theorem}

\begin{proof}[Sketch]
Stress-energy conservation across the transition requires the phase-C vacuum
energy to land somewhere. The cascade's metric content after $N = 217$ has
$w_{\mathrm{eff}} = 1/3$ (Paper~III~\S14.6: the Lorentzian cascade scale factor
$a(t) = \sqrt{1 - t^2}$ gives $\rho_{\mathrm{eff}} \propto 3/a^4$,
$p_{\mathrm{eff}} \propto 1/a^4$, hence $w_{\mathrm{eff}} = 1/3$). This is the
only $w = 1/3$ content structurally available in the cascade at $d = 4$. The
released vacuum energy therefore appears as radiation in the Friedmann
equation.
\end{proof}

\subsection{Reheating temperature}

Applying $\rho_r = (\pi^2/30)\,g_{\mathrm{eff}}\,T^4$ to the released energy:

\begin{corollary}[Reheating temperature]\label{cor:TRH}
\[
T_{\mathrm{RH}} \;=\; \left(\frac{30\,\Delta\rho}{\pi^2\,g_{\mathrm{eff}}}\right)^{1/4}
M_{\mathrm{Pl,red}}.
\]
Using $\Delta\rho = 0.30\,M_{\mathrm{Pl,red}}^4$ and the cascade-native
$g_{\mathrm{eff}} = 106.75$ derived below:
\[
T_{\mathrm{RH}} \;=\; 0.304\,M_{\mathrm{Pl,red}} \;=\; 7.40 \times 10^{17}\,\mathrm{GeV}.
\]
\end{corollary}

\subsection{Cascade-native $g_{\mathrm{eff}}$}\label{subsec:g_eff}

The counting of relativistic degrees of freedom at $T_{\mathrm{RH}}$ is a
cascade-internal calculation, not an SM import. Each distinguished layer
below the phase transition $d_1 = 19$ (where fermion projections are
subcritical; Part~IVa \S4), plus the gauge window $\{12, 13, 14\}$ and
the Gen~1 layer at $d = 21$ (subcritical through the transition width,
Part~IVa Thm~4.4), contributes its 4D-projected propagating modes.
Following Part~IVa~\cite{paper4a} \S4 ``The cascade obstruction map'':

\begin{center}
\renewcommand{\arraystretch}{1.25}
\begin{tabular}{ccccc}
\toprule
$d$ & Cascade content & Bosons & Fermions & Role \\
\midrule
$5$  & Gen~3 Dirac                         & $0$   & $30$ & $\tau$, $b$, $\nu_\tau$ \\
$12$ & $\mathrm{SU}(3)$ gauge              & $16$  & $0$  & $8$ gluons $\times\,2$ pol. \\
$13$ & $\mathrm{SU}(2)$ + Gen~2 + Higgs    & $10$  & $30$ & $3\,W + $ Higgs $+$ $\mu,s,c$ \\
$14$ & $\mathrm{U}(1)$ gauge               & $2$   & $0$  & $B$ boson $\times\,2$ pol. \\
$21$ & Gen~1 Dirac                         & $0$   & $30$ & $e$, $d$, $u$ \\
\midrule
\multicolumn{2}{l}{\textbf{Total}}          & $28$  & $90$ & \\
\bottomrule
\end{tabular}
\end{center}

Per-generation fermion count $30$: $1$ charged lepton (Dirac, 4 d.o.f.) +
$1$ left-handed neutrino (2 d.o.f.) + $1$ up-type quark with 3 colours
($3 \times 4 = 12$ d.o.f.) + $1$ down-type quark with 3 colours
($12$ d.o.f.). Higgs content $4$: complex doublet in the unbroken phase
(high-$T$, above the EW symmetry-breaking scale, which is automatic at
$T \gg 10^{2}$~GeV~$\ll T_{\mathrm{RH}}$). Gauge bosons count $2$
polarisations each at high-$T$ (massless unbroken phase).

\begin{proposition}[Cascade-native $g_{\mathrm{eff}}$]\label{prop:g_eff}
Summing over distinguished cascade layers, with fermion d.o.f.\ weighted by
$7/8$:
\[
g_{\mathrm{eff}}^{\mathrm{cascade}}
\;=\; 28 \;+\; \frac{7}{8}\cdot 90 \;=\; 106.75.
\]
This equals the Standard Model value exactly. The equality is not an
assumption: Part~IVa~\cite{paper4a} forces the gauge group
$\mathrm{SU}(3)\times\mathrm{SU}(2)\times\mathrm{U}(1)$ (Adams' theorem +
Bott window), three generations (Bott period + $d_1 = 19$ cutoff), and
one Higgs doublet (unique hairy-ball zero on $S^{12}$). The cascade's
4D projection reproduces the Standard Model spectrum layer-by-layer; the
$g_{\mathrm{eff}}$ count follows directly.
\end{proposition}

\begin{remark}[The $d = 29$ open question]\label{rem:d29}
Part~IVa's Bott period~$8$ places a 4th fermion layer at $d = 29$,
suppressed below $d_1 = 19$ to $\sim\! 0.5$~eV (Cover sheet; Part~IVb
Open Question~5). At $T_{\mathrm{RH}} \sim 10^{18}$~GeV all such masses
are irrelevant: the layer's modes, if they exist as RH neutrino
partners, are relativistic and contribute to $g_{\mathrm{eff}}$. Under
the alternative scenario with three Majorana RH neutrinos (2~d.o.f.\
each) sourced at $d = 29$:
\[
g_{\mathrm{eff}}^{\mathrm{cascade},\,d{=}29}
\;=\; 28 \;+\; \frac{7}{8}\cdot(90 + 6) \;=\; 112.00,
\]
giving $T_{\mathrm{RH}} = 0.3004\,M_{\mathrm{Pl,red}} = 7.31 \times 10^{17}$~GeV
---a $-1.2\%$ shift relative to the cascade-${\equiv}$SM value. The shift
is numerically negligible at the Planck-scale level; it does not resolve
the gravitino-bound exposure (see Section~\ref{sec:falsifiers}).
\end{remark}

\begin{remark}[What this replaces]\label{rem:replaces-SM}
Earlier drafts of this paper carried $g_{\mathrm{eff}} = 106.75$ as a
``Standard Model approximation''---a provisional input flagged as the
first Open Question. Proposition~\ref{prop:g_eff} removes the
provisionality: the SM spectrum is the cascade's own 4D projection, not
an external import. The Tier classification for
Corollary~\ref{cor:TRH} and its input $g_{\mathrm{eff}}$ is correspondingly
updated (Section~\ref{sec:tiers}). The computation is reproducible via
\texttt{tools/generators/cascade\_g\_eff.py}.
\end{remark}

\subsection{Why the Big Bang is not a singularity}

The $N = 217$ transition has:
\begin{itemize}[nosep]
\item \emph{Finite energy density} before and after ($0.30$ and $7 \times 10^{-121}$ in Planck units).
\item \emph{Finite Hubble rate} before ($H / M_{\mathrm{Pl,red}} \approx 0.32$)
and immediately after (inherited from the sum of residual $\rho_\Lambda^{(D)}$
and the newly-thermalised $\rho_r$; simulator output gives $H \sim 0.3$ dropping
to $\sim 0.05$ over the first 30 ticks as radiation dilutes).
\item \emph{Finite scale factor}, continuous across the transition (the
single-tick layer addition does not discontinuously rescale $a$).
\item \emph{No curvature divergence}, no geodesic incompleteness, no breakdown
of the metric.
\end{itemize}

The ``Big Bang'' in this reading is the moment the cascade's second threshold
first exists as a structural feature of the tower. Before $N = 217$ the tower
has not reached the terminus of Paper~0; after $N = 217$ the terminus is
permanent and $\rho_\Lambda$ sits at its asymptotic value. What was a
large-magnitude vacuum energy becomes a small-magnitude vacuum energy plus a
thermal bath, in one tick. No initial singularity is required.

\begin{remark}[Connection to Paper~V \S5.8]
Paper~V~\cite{paper5}, Section~5.8 asserts that the cascade has ``no
singularity, no inflation, no free initial conditions'' because the descent
from $d = \infty$ is a static logical structure rather than a temporal process.
Under Part~VI's tower-growth reading, the ``no singularity'' claim is preserved
(this section), inflation is provided (Section~\ref{sec:phases}), and the
initial conditions are derived from Paper~0's threshold structure
(Theorem~\ref{thm:reheating}). The Paper~V position is refined, not
contradicted: Part~VI supplies the dynamical history that Paper~V's atemporal
framing did not require but can now accommodate.
\end{remark}

\section{Perturbations: The Gram Deficit}
\label{sec:perturbations}

\subsection{The cascade's native source of inhomogeneity}

The cascade has no inflaton scalar field and no quantised metric operator, so
the standard sources of cosmological perturbation are not available. What the
cascade does have is the layer-to-layer overlap deficit of the slicing
integrands. At adjacent cascade layers $d$ and $d+1$, the integrands
$f_d(x) = (1-x^2)^{d/2}$ are not collinear in $L^2$. Their normalised overlap
satisfies (Paper~0 Supplement~\cite{paper0supp} Theorem~15.2):
\[
1 - C_{d,d+1}^2 \;=\; 1 \,-\, \frac{B\!\left(\frac{1}{2},\, d + \frac{3}{2}\right)^2}
{B\!\left(\frac{1}{2},\, d + 1\right)\,B\!\left(\frac{1}{2},\, d + 2\right)}
\;\sim\; \frac{1}{8\,d^2}
\quad (d \to \infty).
\]
This is a Beta-function identity with no free parameters. It is the cascade's
native ``vacuum fluctuation'' structure: each tick adds a layer whose integrand
is almost but not exactly aligned with the one below it, and the non-aligned
fraction is the cascade-internal perturbation seeded at that tick.

\begin{definition}[Cascade-native perturbation amplitude]\label{def:gram-amp}
At layer $d$, added to the tower during inflation, the seeded perturbation
amplitude is
\[
\delta_d \;=\; \sqrt{1 - C_{d,d+1}^2}
\;\sim\; \frac{1}{2\sqrt{2}\,d} \qquad (d \to \infty).
\]
The cascade's natural index is the layer dimension $d$, not a comoving
wavenumber $k$.
\end{definition}

\subsection{Spectrum in the cascade's natural index}

Evaluating $\delta_d$ across the pre-Big-Bang phases, using the exact
Beta-function formula (not the asymptotic $1/(2\sqrt{2}\,d)$ of
Definition~\ref{def:gram-amp}):

\begin{center}
\renewcommand{\arraystretch}{1.25}
\begin{tabular}{llcc}
\toprule
Phase & $d$ range & $\delta_d$ range (exact) & Asymptotic $1/(2\sqrt{2}\,d)$ \\
\midrule
A & $4$--$6$ & $0.067$--$0.049$ & $0.088$--$0.059$ \\
B & $7$--$18$ & $0.043$--$0.018$ & $0.051$--$0.020$ \\
C & $19$--$216$ & $0.018$--$0.0016$ & $0.019$--$0.0016$ \\
\bottomrule
\end{tabular}
\end{center}

Values in the ``exact'' column come from
$\delta_d = \sqrt{1 - C_{d,d+1}^2}$ with
$C_{d,d+1} = B(\tfrac{1}{2},\,d+\tfrac{3}{2}) / \sqrt{B(\tfrac{1}{2},\,d+1)\,B(\tfrac{1}{2},\,d+2)}$
evaluated numerically. The asymptotic $1/(2\sqrt{2}\,d)$ is reliable within
$\sim 10\%$ for $d \gtrsim 19$ (Phase~C) but overestimates by $20$--$40\%$
in Phase~A, where finite-$d$ corrections to the Beta-function ratio are
non-negligible. The asymptotic form remains useful for scaling arguments
and for the conversion-to-$k$-space derivation below; all numerical
comparisons to observation should use the exact values.

\begin{proposition}[Cascade power spectrum in layer index]\label{prop:dspectrum}
The cascade's perturbation power spectrum, indexed by the layer dimension $d$
at which each mode is seeded, is
\[
P_\delta(d) \;=\; |\delta_d|^2 \;\sim\; \frac{1}{8\,d^2}.
\]
The spectrum is monotonically decreasing in $d$: modes seeded deeper in the
tower (larger $d$) carry smaller amplitudes. No free parameters enter.
\end{proposition}

\begin{remark}
Paper~IVb's $\alpha(d^*)/\chi^k$ family (Part~IVb~\cite{paper4b} Remark~4.6)
uses the quantity $\alpha(d) = R(d)^2/4$ at distinguished cascade layers. The
Gram amplitude $\delta_d \sim 1/(2\sqrt{2}\,d)$ scales the same way as
$\sqrt{\alpha(d)}$ does asymptotically ($R(d) \sim \sqrt{2/d}$ for large $d$).
This is not coincidence: both derive from the Beta-function asymptotics of the
slicing recurrence. The Gram deficit and the cascade coupling are two views of
the same per-layer residual.
\end{remark}

\subsection{Conversion to comoving wavenumber}

Observational cosmology parametrises perturbations by comoving wavenumber $k$,
not by cascade layer $d$. The conversion under tower-growth is set by the
inflation history.

During Phase~C, $H / M_{\mathrm{Pl,red}} \approx 0.317$ is approximately
constant. With $\Delta t = t_{\mathrm{Pl,red}}$ per tick and
$a(t) \propto \exp(H\,t)$, the scale factor at the moment layer $d$ is added
satisfies
\[
\ln\!\left(\frac{a(d)}{a(19)}\right) \;=\; H \cdot (d - 19) \cdot t_{\mathrm{Pl,red}}
\;=\; 0.317\,(d - 19).
\]
The comoving scale exiting the horizon at tick $d$ is $k_d = a(d)\,H$, so
\[
\ln k_d \;=\; \mathrm{const} \,+\, 0.317\,(d - 19),
\]
linear in $d$ within Phase~C. Inverting: $d \propto \ln k$.

\begin{corollary}[Spectrum in $k$-space]\label{cor:kspectrum}
Under tower-growth inflation, the cascade-native perturbation amplitude seeded
at comoving wavenumber $k$ satisfies
\[
\delta(k) \;\propto\; \frac{1}{\ln k} \quad\text{within Phase~C.}
\]
Equivalently, the power spectrum is
\[
P_\delta(k) \;\propto\; \frac{1}{(\ln k)^2} \quad\text{within Phase~C.}
\]
\end{corollary}

\subsection{Comparison to observed $n_s$}

The observed CMB scalar power spectrum is near-power-law:
\[
P_s(k) \;\propto\; k^{n_s - 1}, \qquad n_s = 0.9649 \pm 0.0042\ \text{(Planck 2018).}
\]
Expanding a power law: $k^{n_s - 1} = 1 + (n_s - 1)\ln k + \tfrac{1}{2}(n_s - 1)^2(\ln k)^2 + \cdots$. Expanding the cascade's $1/(\ln k)^2$ around a reference scale
$k_0$: $1/(\ln k)^2 = 1 - 2\,\delta\!\ln k / \ln k_0 + \cdots$ where
$\delta\!\ln k = \ln(k/k_0)$. Within Planck's observational range
($k \sim 10^{-4}$ to $10^{-1}\,\mathrm{Mpc}^{-1}$), both expressions admit
near-logarithmic behaviour at leading order. At Planck precision the two forms
are not cleanly distinguishable.

\begin{proposition}[Running of $n_s$]\label{prop:running}
The cascade's $P_\delta(k) \propto 1/(\ln k)^2$ form predicts a non-zero
running $\alpha_s \equiv d n_s / d\!\ln k$ with a specific functional structure:
\[
\alpha_s^{\mathrm{cas}}(k) \;\simeq\; \frac{-2}{(\ln k)^2}
\]
at leading order in the expansion. A pure power law has $\alpha_s = 0$ exactly;
the cascade predicts $\alpha_s$ is non-zero and scale-dependent.
\end{proposition}

\begin{remark}[Current bound]
Planck 2018 gives $\alpha_s = -0.0045 \pm 0.0067$, consistent with zero at
$1\sigma$. The cascade's structural prediction is of the same order but with
specific $k$-dependence; CMB-S4 aims for $\sigma(\alpha_s) \sim 0.002$, which
would begin to probe the cascade's functional form.
\end{remark}

\subsection{Falsifier}

\begin{proposition}[CMB-S4 falsifier]\label{prop:CMBS4-falsifier}
If CMB-S4 measures $\alpha_s$ consistent with zero to $\lesssim 10^{-3}$
precision, the tower-growth perturbation spectrum of
Corollary~\ref{cor:kspectrum} is disfavoured: a pure power-law $k^{n_s-1}$
with $\alpha_s = 0$ exactly is incompatible with the cascade's
$P_\delta \propto 1/(\ln k)^2$.
\end{proposition}

\begin{remark}[What this falsifier does not do]
This falsifier tests the specific tower-growth reading of Part~VI, not the
cascade series generally. Parts~0--V do not predict a perturbation spectrum
and are not challenged by any $n_s$ or $\alpha_s$ measurement.
\end{remark}

\begin{remark}[Perturbation freeze-out]
Standard inflation freezes quantum perturbations at horizon exit, converting
them to classical density perturbations. The cascade has no horizon exit in
the standard sense (the tower grows discretely; there is no continuous
horizon-crossing event per mode). Whether the Gram-deficit seeds become
classical via geometric decoherence (Paper~II~\cite{paper2} Section~8.5) acting
over subsequent layers, or via a different mechanism, is open
(Section~\ref{sec:open}, Open~Question~4). The spectrum of
Proposition~\ref{prop:dspectrum} is the amplitude \emph{at seed}; its late-time
evolution under post-Big-Bang dynamics requires the freeze-out mechanism to be
specified.
\end{remark}

\section{Tensor Modes}
\label{sec:tensors}

\subsection{No cascade analogue of $r = 16\epsilon$}

Standard slow-roll inflation derives the tensor-to-scalar ratio as
$r = P_t / P_s = 16\,\epsilon$, where $P_t$ is the power spectrum of quantum
fluctuations of the metric operator $\hat{g}_{\mu\nu}$ at horizon exit and
$\epsilon$ is the slow-roll parameter of the inflaton potential. The cascade
has neither ingredient:

\begin{itemize}[nosep]
\item The metric is a state property of $|\Psi\rangle \in S^{d-1}$
(Paper~II\,=\,III~\cite{paper23} Theorem~6.1), not a quantum operator on
a fixed background. The metric is never promoted to $\hat{g}_{\mu\nu}$;
fluctuations of $\hat{g}$ are not part of the cascade's mathematical
structure.
\item There is no inflaton scalar field. The ``inflaton'' during Phase~C is
the cascade's discrete phase-dependent cosmological constant, which takes a
single value throughout the phase. There is no rolling potential and no
slow-roll parameter.
\end{itemize}

Both ingredients of the $r = 16\epsilon$ derivation are therefore absent. The
formula has no cascade analogue. Whatever tensor signal the cascade predicts
is \emph{not} computed via this route.

\subsection{What tensor modes are in the cascade}

\begin{proposition}[Tensor modes as classical induced responses]\label{prop:tensor-induced}
Under Part~VI's tower-growth reading, tensor perturbations at the observer's
frame $d = 4$ are classical responses of the projected metric to the same
Gram-deficit scalar source of Section~\ref{sec:perturbations}, via the
Einstein equation of Paper~III~\cite{paper3}. They are not independent
quantum degrees of freedom.
\end{proposition}

\begin{proof}[Sketch]
The cascade has one quantum object: the cascade state $|\Psi\rangle$ on
$S^{d-1}$. Its perturbations, sourced by the Gram deficit between consecutive
cascade layers, enter the 4D effective stress-energy tensor via Lovelock's
projection identity (Paper~III~\S8). The Einstein equation
$G_{\mu\nu} + \Lambda g_{\mu\nu} = 8\pi G\,T_{\mu\nu}$ is the unique
gravitational equation at $d = 4$ (Paper~III Theorem~3.1), so the metric
response to a stress-energy perturbation is uniquely determined. In
particular, any spin-2 component of the induced metric response is a
\emph{classical} derivation from the scalar source, not a separately-quantised
field. No spin-2 quantum field exists in the cascade; the induced classical
tensor is the sole spin-2 content at $d = 4$.
\end{proof}

\begin{remark}[Relation to Paper~III \S12]
Paper~III~\S12 already commits to ``no gravitons'' on the grounds that the
metric is a state property rather than an operator. Proposition~\ref{prop:tensor-induced}
applies that commitment to primordial cosmology: there is no graviton vacuum
fluctuation to source a primordial tensor background; only the classical
induced response of the projected metric remains.
\end{remark}

\subsection{Magnitude of the induced tensor response: open}

Computing the magnitude of the cascade-native tensor-to-scalar ratio requires:

\begin{enumerate}[nosep]
\item Translating the Gram-deficit scalar source (Definition~\ref{def:gram-amp})
into a perturbation $\delta T_{\mu\nu}$ of the effective stress-energy tensor
at $d = 4$.
\item Solving the linearised Einstein equation $\delta G_{\mu\nu} = 8\pi G\,\delta T_{\mu\nu}$ for the induced metric response $\delta g_{\mu\nu}$.
\item Decomposing $\delta g_{\mu\nu}$ into scalar, vector, and tensor parts
under $\mathrm{SO}(3)$ on the observer's spatial $S^3$, and propagating the
tensor part through Phase~D expansion to recombination.
\item Reading off the tensor-to-scalar amplitude ratio at CMB scales.
\end{enumerate}

None of these steps has been performed rigorously. The author's preliminary
discussion (draft notes, \emph{not} in the published series) offered a
hand-wave estimate $r \sim e^{-60}$ via naive $e^{-N/2}$ suppression of
super-horizon tensor modes through Phase~D expansion. This reasoning is
incorrect: super-horizon tensor modes freeze rather than decay, so no such
exponential suppression operates. The preliminary estimate is retracted. The
correct calculation has not been performed.

\begin{proposition}[Provisional tensor suppression]\label{prop:tensor-provisional}
The cascade's classical-induced-tensor mechanism produces no primordial
gravitational-wave background of the quantum-inflaton-fluctuation kind. The
amplitude of the classical induced tensor response is structurally suppressed
relative to the scalar source, but the magnitude of the suppression is not
currently computable from the cascade's own structure without further
technical work.
\end{proposition}

\begin{remark}[Why this matters for falsification]
Proposition~\ref{prop:tensor-provisional} is weaker than a sharp prediction of
$r$. It commits to ``no standard primordial gravitational wave signal'' but
does not commit to a numerical value of the residual induced-response
amplitude. A positive detection at any $r$ detectable by near-term experiments
forces one of two conclusions: either the cascade's ``no graviton'' commitment
is wrong (falsifying Paper~III~\S12 and Paper~II\,=\,III Theorem~6.1), or the
induced classical response happens to be non-negligible and the cascade's
rigorous $r$ calculation, once performed, reproduces the detection.
\end{remark}

\subsection{Falsifier}

\begin{proposition}[LiteBIRD falsifier]\label{prop:litebird}
A detection of primordial B-mode polarisation at $r \gtrsim 10^{-3}$ with
spectral shape characteristic of inflaton-quantum-fluctuation origin
(tensor tilt $n_t \approx -r/8$ from slow-roll consistency) falsifies the
cascade's ``no graviton'' commitment and therefore the tower-growth reading.
Non-detection at LiteBIRD sensitivity is compatible with the cascade but does
not confirm it; the cascade's induced-response mechanism may produce a signal
at any magnitude from ``unobservable'' to ``near the LiteBIRD floor'',
pending the rigorous calculation.
\end{proposition}

\begin{remark}
The cleanest form of the LiteBIRD falsifier is not the amplitude alone but
the \emph{spectral shape}. A slow-roll inflaton prediction has $n_t < 0$ set
by the consistency relation; a cascade-induced classical response has no
reason to satisfy that relation and would exhibit different $n_t$ behaviour.
If a future tensor detection includes spectral information, the slow-roll
consistency relation becomes a pointed test of the cascade's alternative
origin.
\end{remark}

\section{Particle Structural Timeline}
\label{sec:particles}

\subsection{Layer addition and topological features}

Each cascade layer, when added to the truncated tower, brings its topological
features into structural existence. Parts~IVa and~IVb~\cite{paper4a,paper4b}
identify what each distinguished layer carries: Dirac obstructions at fermion
layers, Adams vector fields at gauge layers, hairy-ball zeros at
even-dimensional spheres. Under tower-growth, the \emph{time} at which each
feature becomes structurally available is the tick at which its layer is added.

\begin{center}
\renewcommand{\arraystretch}{1.25}
\begin{tabular}{cccl}
\toprule
Added at tick & Layer & Feature & Source \\
\midrule
$N = 5$   & $d = 5$   & Gen~3 Dirac obstruction ($\tau$, $b$, $\nu_\tau$); observer host & Paper~IVa~\cite{paper4a} \S4 \\
$N = 7$   & $d = 7$   & Area maximum $d_0$; decay-rate zero & Paper~0~\cite{paper0} \S5 \\
$N = 12$  & $d = 12$  & $\mathrm{SU}(3)$ layer (Adams $\rho(12)-1=3$) & Paper~IVa \S2.5 \\
$N = 13$  & $d = 13$  & $\mathrm{SU}(2)$ layer + Gen~2 Dirac obstruction ($\mu$, $s$, $c$) & Paper~IVa \S2--\S4 \\
$N = 14$  & $d = 14$  & $\mathrm{U}(1)$ layer & Paper~IVa \S2 \\
$N = 19$  & $d = 19$  & Phase-transition threshold $d_1$ & Paper~0 Thm~6.7 \\
$N = 21$  & $d = 21$  & Gen~1 Dirac obstruction ($e$, $d$, $u$) & Paper~IVa \S4 \\
$N = 217$ & $d = 217$ & Second threshold $d_2$; Big Bang & Paper~0 Thm~6.7 \\
\bottomrule
\end{tabular}
\end{center}

The ordering is strict and cascade-internal: colour ($d = 12$) precedes weak
isospin ($d = 13$) by one tick, which precedes hypercharge ($d = 14$) by one
tick. Fermion generations are spaced by the Bott period: Gen~3 ($d = 5$) to
Gen~2 ($d = 13$) to Gen~1 ($d = 21$), eight ticks apart. The Standard Model's
structural features assemble themselves during Phase~C of inflation in a
specific order that is not adjustable.

\subsection{Structural existence versus particle occupancy}

\begin{remark}[What ``structural existence'' means]
A layer is \emph{structurally present} at time $t$ if it is part of the
truncated tower $[N_0, N(t)]$. Structural presence implies: the layer's
sphere-area contributes to cascade sums; its topological obstruction (hairy-ball
zero if applicable) exists as part of the active tower topology; its gauge or
fermion projection is available as a feature that any cascade state can couple
to.

It does \emph{not} imply that particles with the layer's quantum numbers
physically occupy the layer during the pre-Big-Bang phase. Pre-$N = 217$ the
universe has zero matter density and zero radiation density; only vacuum energy
exists (Section~\ref{sec:phases}). The ``particles'' of the Standard Model are
features of the \emph{post}-$N = 217$ tower projected to $d = 4$ under
post-Big-Bang dynamics; their structural layers are added during inflation, but
their observable occupancy does not begin until the $N = 217$ thermalisation
produces a radiation bath.
\end{remark}

\subsection{Structural prediction}

\begin{proposition}[Cascade-internal ordering]\label{prop:ordering}
Any cosmological observable that distinguishes Standard Model features by
cascade-internal addition order must respect the ordering of the table above.
Specifically:
\begin{itemize}[nosep]
\item $\mathrm{SU}(3)$ structural availability precedes $\mathrm{SU}(2)$ and
$\mathrm{U}(1)$ by one and two ticks respectively.
\item Gen~3 fermions (heavier) are structurally available before Gen~2, which
precedes Gen~1 (lightest) --- the opposite of the ``heavy particles froze out
first'' story of standard cosmology, which is a post-Big-Bang thermal
statement about occupancy at a given temperature.
\item Gen~1 fermions ($e$, $d$, $u$), the lightest generation and the observer's
own matter content, become structurally available last (Phase~C, at $N = 21$).
\end{itemize}
\end{proposition}

\begin{remark}[No conflict with thermal freeze-out]
Proposition~\ref{prop:ordering} is not a claim about particle freeze-out
temperatures or abundance ratios. Standard cosmology's ``heavy particles
froze out first'' ordering is a statement about post-Big-Bang thermal
equilibrium and decoupling at specific temperatures $T \sim m$. The
cascade-internal ordering is about \emph{structural availability} during the
pre-Big-Bang tower-growth phase. Both orderings are compatible: the structural
features exist in the cascade order during inflation; the observable
occupancies then freeze out in thermal order after the Big Bang according to
standard Boltzmann-equation dynamics with cascade-derived masses.
\end{remark}

\begin{remark}[A potential observable signature]
If some post-Big-Bang observable retains memory of the \emph{order} in which
cascade features became structurally available --- distinct from the order in
which particle occupancies froze out --- this would be a cascade-native
signature. Candidates include non-thermal relic abundances set by
layer-addition rather than by thermal decoupling, or specific non-Gaussian
correlations in the CMB between modes seeded at different cascade ticks. No
concrete such signature has been identified; this is flagged as
Section~\ref{sec:open} Open~Question~9.
\end{remark}

\section{Retroactive Clarifications to Paper~II}
\label{sec:paper2-amendments}

Part~VI's dynamical reading clarifies three points in Paper~II~\cite{paper2}
without contradicting any of its theorems. Each clarification is a refinement
under the tower-growth identification; the within-moment content of Paper~II
is preserved unchanged.

\subsection{Paper~II \S7.1: the time identification}

Paper~II~\S7.1 states: ``The cascade does not happen in time; the cascade is
time. The arrow of time is the direction of descent, which is irreversible.''

\begin{remark}[Clarification]
This remains correct. Part~VI refines it by identifying \emph{two} coupled
operations both called ``time'' in Paper~II:
\begin{itemize}[nosep]
\item the \emph{within-moment descent} from $d = N(t)$ down to $d = 4$ through
the currently-existing tower (Paper~II's primary reading), and
\item the \emph{between-moments growth} of the tower, $N(t) \to N(t) + 1$,
which Paper~II does not distinguish.
\end{itemize}
Paper~II's ``arrow of time = direction of descent'' captures the first.
Physical time evolution between moments is the second. The irreversibility of
Paper~II \S7 (slicing compresses one direction, lost information resides on
the boundary via $\Omega_{d-1}/V_d = d$) applies to the within-moment
projection. The between-moments layer-addition is also irreversible, on
different grounds: once a layer is added at the top of the tower, the tower
cannot shrink without violating cascade-content-preservation across the
transition.
\end{remark}

\subsection{Paper~II \S7.3: the discrete propagator}

Paper~II~\S7.3 Theorem~7.1 states that $K(D, d') = \prod_{j = d'}^{D-1} L(j)$
is the cascade propagator, with $L(j) = iN(j)$ in the precessing cascade.

\begin{remark}[Clarification]
Under Part~VI, $K$ is the \emph{within-moment projection} through the
currently-existing tower, not the physical time-evolution operator between
moments. Physical evolution from $t$ to $t + \alpha\,t_{\mathrm{Pl,red}}$
is the composition (Proposition~\ref{prop:time-evolution})
\[
K(N(t) + 1, 4) \;\circ\; A_{N(t)+1} \;\circ\; K(N(t), 4)^{-1},
\]
where $A_{N(t)+1}$ is the tower-growth operator adding layer $N(t)+1$ at the
top. Paper~II's derivation of the effective Schr\"odinger equation
(\S7.4 Corollary~7.4) is the coarse-grained description of the within-moment
$K$ and is correct on that reading; the between-moments $A$ step has no
continuous analogue and does not appear in the effective Schr\"odinger
equation. This answers a latent question raised but not explicitly addressed
by Paper~II: ``why does the cascade's discrete propagator act as a dynamical
time evolution rather than a static projection?'' Under Part~VI, the answer is:
it does not. The dynamics is $A \circ K$; Paper~II's $K$ is the static
projection factor.
\end{remark}

\subsection{Paper~II \S8.6: cumulative decoherence}

Paper~II~\S8.6 reports the eigenvalue deficit of the Gram correlation matrix
over the full 213-layer path $d = 5, \ldots, 217$:
$\varepsilon = 1 - \lambda_1/n = 0.0611$ (the cascade retains $93.9\%$
coherence over the full descent).

\begin{remark}[Clarification]
The $\varepsilon = 6.11\%$ figure assumes the full cascade tower is always
present. Under Part~VI's tower-growth reading, the tower at time $t$ is
$[N_0, N(t)]$, so the Gram matrix at time $t$ is the submatrix over the
currently-existing layers. The time-dependent eigenvalue deficit
\[
\varepsilon(t) \;=\; 1 \,-\, \frac{\lambda_1\bigl(C_{[N_0, N(t)]}\bigr)}{N(t) - N_0 + 1}
\]
grows monotonically with $N(t)$, saturating at $\varepsilon(217) = 6.11\%$
when the full tower is present. Exact values (with $N_0 = 5$, matching
Paper~II's path $d = 5, \ldots, 217$) are computed numerically from the
Beta-function Gram matrix of Paper~0 Supplement~\cite{paper0supp}
Theorem~15.1.
\end{remark}

\begin{center}
\renewcommand{\arraystretch}{1.2}
\begin{tabular}{cccc}
\toprule
Tick $N$ & Layers in tower & $\varepsilon(N)$ & Fraction of $\varepsilon(217)$ \\
\midrule
$7$   & 3   & $0.18\%$ & $3.0\%$ \\
$19$  & 15  & $1.63\%$ & $26.7\%$ \\
$50$  & 46  & $3.70\%$ & $60.5\%$ \\
$100$ & 96  & $5.03\%$ & $82.3\%$ \\
$217$ & 213 & $6.11\%$ & $100\%$ (Paper~II saturated value) \\
\bottomrule
\end{tabular}
\end{center}

Values above are computed from the exact Gram matrix using
\texttt{tools/model\_checks/cascade\_decoherence\_vs\_tower\_height.py}; no approximation
enters beyond floating-point arithmetic on the Beta-function entries. The
dependence on $N_0$ is weak: for $N_0 = 4$ the saturated value at $N = 217$
is $\varepsilon = 6.33\%$ (one additional layer at the bottom contributes
$\sim 0.2\%$), with all intermediate values correspondingly $\sim 0.1\%$--$1\%$
larger.

\begin{corollary}[Pre-recombination Gram correction]\label{cor:pre-rec-gram}
Any descent-dependent cascade observable set at an epoch before $N = 217$
should use $\varepsilon(N_{\mathrm{set}})$ rather than $\varepsilon(217)$ in
the Part~0 Supplement first-order correction formula. For observables set
during Phase~C at $N \sim 100$--$200$, the appropriate correction is
$\varepsilon(N) \in [5.0\%, 6.0\%]$, slightly smaller than the saturated
value. The numerical impact on Part~V~\cite{paper5} observables is
negligible --- all Part~V observables (cosmological parameters, $H_0$, $r_d$,
$t_0$) are set at post-Big-Bang epochs where $\varepsilon(N) =
\varepsilon(217) = 6.11\%$ to the precision required.
Corollary~\ref{cor:pre-rec-gram} is a structural refinement that becomes
numerically relevant only for future cascade observables set at specific
pre-recombination epochs.
\end{corollary}

\begin{remark}[Does not modify Part~V]
Paper~V's Gram-corrected predictions for $\rho_\Lambda$, $H_0$, $T_{\rm CMB}$,
$\Omega_m$, etc., use the saturated $\varepsilon = 6.11\%$. Under Part~VI's
refinement, this remains the correct value for all Part~V observables because
all of them are set at post-Big-Bang epochs. No numerical adjustment to
Part~V is required. Part~VI's clarification is a structural statement about
\emph{when} the saturated value applies, sharpening but not changing the
current numerics.
\end{remark}

\section{Connection to the Cover-Sheet Thought Experiment}
\label{sec:coversheet}

The cover sheet describes the observer on the $S^3$ shell of a 5D black hole,
asymptotically falling toward the pseudo-horizon. Tower growth is the dual of
this picture:

\begin{center}
\renewcommand{\arraystretch}{1.25}
\begin{tabular}{ll}
\toprule
Cover-sheet (infall) & Tower growth (this paper) \\
\midrule
Shell falls toward horizon & Tower grows at top \\
Asymptotic approach, never completes & $N \to \infty$, never completes \\
One step deeper along infall axis & One layer added above observer \\
Pseudo-horizon recedes as BH evaporates & Terminus at $N = 217$ permanent post-Big-Bang \\
Observer on $S^3$ shell of $d = 5$ BH & Observer at $d = 4$, host at $d = 5$ (\S\ref{subsec:N0}) \\
Infinite time dilation at the horizon & Infinite layers above observer at $N \to \infty$ \\
\bottomrule
\end{tabular}
\end{center}

Both describe the same asymptotic process in the cascade's geometry. The
cover-sheet picture frames the observer as falling; Part~VI frames the tower
as growing. Neither process completes in finite time; both saturate
asymptotically.

\begin{remark}[Locking the time identification]
The cover-sheet identification ``time is the observer's asymptotic infall
toward the pseudo-horizon'' is used by Paper~I~\cite{paper1} Section~3.2 to
justify the locked-time 3D projection factor $\Omega_2/\Omega_3 = 2/\pi$.
Under tower growth, this identification becomes: time is the asymptotic
growth of the cascade tower above the observer. The two formulations agree on
the key feature that forces the Paper~I factor --- time is not a free fourth
dimension --- and agree on the asymptotic character that makes the descent
never complete. Part~VI does not alter Paper~I's derivation; it supplies an
equivalent dynamical reading of the ``locked time'' hypothesis.
\end{remark}

\begin{remark}[Three unified pictures of time]
The cascade series has now accumulated three descriptions of time, all
mutually consistent under Part~VI's identification:
\begin{enumerate}[nosep]
\item \emph{Paper~II \S7}: time is the cascade's slicing direction; descent
from high $d$ to low $d$ is the arrow of time. (Within-moment reading.)
\item \emph{Cover sheet}: time is the observer's asymptotic infall toward the
5D-BH pseudo-horizon. (Geometric dual.)
\item \emph{Part~VI}: time is the growth of the truncated tower above the
observer; one Planck tick adds one layer. (Dynamical reading.)
\end{enumerate}
These are three vocabularies for one physical fact: the cascade's structural
descent is the observer's experience of time. The three agree on the key
properties (asymptotic, never-completing, locked, irreversible) and differ
only in which operation is primary. Part~VI takes (3) as primary and derives
the cosmic-history consequences.
\end{remark}

\section{Falsifiers}
\label{sec:falsifiers}

Part~VI is explicitly speculative and makes several sharp predictions that
near-future experiments can test. The falsifiers listed below apply to the
tower-growth reading specifically; none challenges Parts~0--V.

\begin{enumerate}[leftmargin=*,nosep]

\item \textbf{CMB-S4 running of $n_s$.} The cascade's $P_\delta(k) \propto
1/(\ln k)^2$ spectrum predicts a specific running $\alpha_s \equiv
d n_s / d\!\ln k \sim -2/(\ln k)^2$ (Proposition~\ref{prop:running}). A pure
power law has $\alpha_s = 0$ exactly. If CMB-S4 measures $\alpha_s$ consistent
with zero at $\lesssim 10^{-3}$ precision, the tower-growth spectrum is
disfavoured.

\item \textbf{LiteBIRD tensor detection with slow-roll consistency shape.} A
primordial B-mode signal at $r \gtrsim 10^{-3}$ with tensor tilt $n_t \approx
-r/8$ (slow-roll consistency relation) falsifies the no-graviton commitment
(Proposition~\ref{prop:litebird}). Non-detection is compatible but not
confirmatory.

\item \textbf{Total e-folds outside $[65, 82]$.} The two cascade-native tick
rates give $N_e = 70.94$ (reading (i), $\alpha = 1$) and $N_e = 81.6$
(reading (ii), $\alpha = 3\pi/8$). If CMB-S4 tightens $N_e$ outside this
combined window, both readings are falsified.

\item \textbf{Reheating temperature.} Gravitino abundance constraints in
supergravity-compatible scenarios and CMB tensor bounds jointly impose
$T_{\mathrm{RH}} \lesssim 10^{15}$--$10^{16}$~GeV in many extensions of
standard cosmology. If these constraints are confirmed at that level, the
cascade's $T_{\mathrm{RH}} \sim 7 \times 10^{17}$~GeV from
Corollary~\ref{cor:TRH} is challenged --- though the bounds themselves rely
on particle content beyond the cascade's own framework.

\item \textbf{Matter density before $N = 217$.} The cascade predicts zero
matter and zero radiation during Phases~A--C (only vacuum energy exists). Any
observational evidence of matter content sourced before the cascade-native
Big Bang (e.g., primordial black holes formed during a pre-$N=217$ matter
era) falsifies the Part~VI account of the pre-Big-Bang universe.

\item \textbf{Discovery of a scalar field driving inflation.} If primordial
non-Gaussianity, B-modes, or cross-correlations between observables jointly
constrain the inflationary dynamics to a specific inflaton potential shape
(at field-theoretic precision), the cascade's no-inflaton reading is
falsified. The cascade commits to inflation being driven by discrete
phase-dependent $\rho_\Lambda$ values, not by a rolling scalar.

\item \textbf{Cascade-order sensitive observables.} If any post-Big-Bang
observable is discovered that distinguishes Standard Model features by
thermal decoupling order but is inconsistent with the cascade-internal
ordering of Proposition~\ref{prop:ordering} (e.g., relic abundance patterns
that favour Gen~1 features over Gen~3 in structural availability), the
particle-timeline prediction is challenged.

\end{enumerate}

The falsifiers above are not uniform in strength. Items~1 and~2 are sharp
and testable at near-future experimental precision. Items~3--5 test
particular cascade-native numerical values that depend on tick-rate and
$g_{\mathrm{eff}}$ choices flagged as open. Items~6 and~7 are qualitative and
require additional work to sharpen.

\section{What This Paper Does and Does Not Do}

\textbf{Does:}
\begin{itemize}[leftmargin=*,nosep]
\item Promotes the truncation height $N$ to a physical time coordinate via
Definition~\ref{def:tower-time}: one Planck tick adds one layer at the top of
the tower, with the observer's dimension $d = 4$ fixed.
\item Splits Paper~II's time identification into within-moment projection
(Paper~II's $K$) and between-moments evolution (tower-growth step $A$);
physical time evolution is the composition $A \circ K$
(Proposition~\ref{prop:time-evolution}).
\item Derives a four-phase pre-Big-Bang structure (Phases~A--D) with
cascade-native $\rho_\Lambda$, $H$, and e-fold count per phase from the
factor-activation at distinguished cascade layers (Section~\ref{sec:phases}).
\item Establishes $N_e = 70.94$ pre-radiation e-folds under tick-rate
reading~(i) (Proposition~\ref{prop:efolds}), inside the observational
inflation window.
\item Identifies the Big Bang as the $N = 217$ discrete phase transition
(Theorem~\ref{thm:reheating}) with reheating temperature $T_{\mathrm{RH}}
\sim 7 \times 10^{17}$~GeV (Corollary~\ref{cor:TRH}) under Standard Model
$g_{\mathrm{eff}}$.
\item Identifies the Gram overlap deficit $\delta_d = \sqrt{1 - C_{d,d+1}^2}$
as the cascade's native source of cosmological perturbations
(Definition~\ref{def:gram-amp}), giving power spectrum $P_\delta(d) \sim
1/(8d^2)$ in the cascade's natural index.
\item Derives the conversion to comoving wavenumber during Phase~C:
$P_\delta(k) \propto 1/(\ln k)^2$ (Corollary~\ref{cor:kspectrum}), predicting
non-zero running $\alpha_s$ with specific $k$-dependence.
\item Argues that $r = 16\epsilon$ has no cascade analogue (neither the
quantised metric operator nor the slow-roll parameter exists); tensor
perturbations at $d = 4$ are classical induced responses of the projected
metric (Proposition~\ref{prop:tensor-induced}).
\item Establishes the cascade-internal particle-structural timeline
(Section~\ref{sec:particles}): Standard Model features assemble in a
specific order as the tower grows during Phase~C.
\item Clarifies three points in Paper~II (\S7.1 time identification, \S7.3
discrete propagator, \S8.6 cumulative decoherence) without contradicting any
of its theorems (Section~\ref{sec:paper2-amendments}).
\item Connects tower-growth to the cover-sheet 5D-BH infall picture as dual
descriptions of the same asymptotic process (Section~\ref{sec:coversheet}).
\item Provides seven specific falsifiers against CMB-S4, LiteBIRD, gravitino
bounds, and structural predictions (Section~\ref{sec:falsifiers}).
\end{itemize}

\textbf{Does not:}
\begin{itemize}[leftmargin=*,nosep]
\item Derive the tick rate $\alpha$ from cascade structure: two candidates
($\alpha = 1$ and $\alpha = 3\pi/8$) are offered, with no structural criterion
to select between them.
\item Derive the initial truncation height $N_0$: $N_0 = 4$ (simulator) and
$N_0 = 5$ (cover-sheet) are both defensible, not yet decided.
\item Rigorously compute the cascade-native tensor-to-scalar ratio: the
magnitude calculation is flagged as explicitly open, and the hand-waved
$r \sim e^{-60}$ estimate from earlier draft discussion is retracted as
incorrect.
\item Derive cascade-native $g_{\mathrm{eff}}$: the Standard Model value
$106.75$ is used provisionally in the reheating calculation.
\item Provide a first-principles mechanism for why layer addition occurs at
one-per-Planck-tick: the identification is proposed, not derived from the
cascade's action.
\item Derive the perturbation freeze-out mechanism: standard slow-roll
horizon-crossing does not apply; geometric decoherence is the candidate
replacement but has not been rigorously shown to convert Gram-deficit quantum
seeds to classical density perturbations at the right amplitude.
\item Produce a full CMB $C_\ell$ prediction: this requires coupling the
Gram-deficit perturbation spectrum to Boltzmann-evolution machinery (CAMB or
CLASS), which has not been done.
\item Modify any Tier 1--3 result of Parts~0--V: Part~VI is Tier~5 throughout
and is explicitly an overlay, not a replacement.
\item Use any result from post-cascade physics: the only inputs are the
cascade's own geometry from Parts~0--V and the tower-growth identification.
\end{itemize}

\section{Open Questions}
\label{sec:open}

\begin{enumerate}[leftmargin=*,nosep]

\item \textbf{Tick rate $\alpha$.} Readings (i) $\alpha = 1$ and (ii) $\alpha
= 3\pi/8$ are both cascade-native; no structural criterion currently
distinguishes them. Option (i) uses the reduced Planck time directly; option
(ii) uses the cascade's lapse at the observer's dimension, $N(4) = 3\pi/8$
(Paper~II \S7.2). CMB-S4 constraints on $N_e$ to $\pm 2$ will distinguish
them observationally; a first-principles derivation would be stronger.

\item \textbf{Initial truncation height $N_0$.} Is $N_0 = 4$ (observer's own
dimension, simulator choice) or $N_0 = 5$ (host dimension, cover-sheet
choice)? The cover-sheet requires the 5D black hole to exist for the
observer's $S^3$ shell to exist, favouring $N_0 = 5$. A structural argument
from Paper~II's boundary dominance ($\Omega_{d-1}/V_d = d$ requires $d \geq 1$
but says nothing about a structural minimum for the observer) has not
resolved the question.

\item \textbf{Cascade-native $g_{\mathrm{eff}}$.}
\emph{Resolved} in Section~\ref{subsec:g_eff}:
Proposition~\ref{prop:g_eff} derives $g_{\mathrm{eff}}^{\mathrm{cascade}} =
106.75$ directly from Part~IVa's layer assignments (gauge window
$\{12, 13, 14\}$, three generations at $d \in \{5, 13, 21\}$, Higgs at
$d = 13$). The equality with the SM value is \emph{forced} by the
cascade's 4D projection, not an assumption. The $d = 29$ open question
(Remark~\ref{rem:d29}) remains: if $d = 29$ RH neutrino partners
thermalise at $T_{\mathrm{RH}}$, the count shifts to $112.00$,
producing a $-1.2\%$ shift in $T_{\mathrm{RH}}$.

\item \textbf{Perturbation freeze-out mechanism.} Standard slow-roll inflation
freezes quantum perturbations at horizon exit. The cascade has no horizon
exit in the standard sense (the tower grows discretely). Candidate mechanism:
geometric decoherence (Paper~II \S8.5) acting over the subsequent layers
between mode seeding and the Big Bang. Whether this produces the correct
classical amplitude at the right epoch is open.

\item \textbf{Tensor-to-scalar ratio $r$, rigorously.} The magnitude of the
classical induced tensor response (Proposition~\ref{prop:tensor-induced})
requires solving the linearised Einstein equation with the Gram-deficit
scalar source and propagating the tensor response through Phase~D. No
rigorous calculation has been performed.

\item \textbf{Pre-$N_0$ ticks.} If $N_0 > 4$, what does the structure at
$N < N_0$ look like? There is no observer (observer requires $d = 4$, needing
$d = 4$ to exist as the bottom of the tower), so the question may be
ill-posed. Whether the cascade admits a ``pre-observer'' phase at all is
open.

\item \textbf{Non-perturbative sector during inflation.} Paper~IVb~\cite{paper4b}
uses $\exp(-\pi/\alpha(5))$ as a non-perturbative factor in the electroweak
VEV $v$. \emph{Partially resolved} by Paper~IVb Remark~4.9: the factor
corresponds to a Kosterlitz--Thouless single-vortex action on the 2-sphere
cross-section of $S^4$ under the Morse foliation of the $d=5$ hairy-ball
defect, not a classical instanton saddle. This suggests a \emph{static}
cascade-geometric mechanism at $d=5$ rather than a dynamical tunnelling
event during inflation, so it does not obviously modify Definition~\ref{def:tower-time}'s
uniform ``one tick per layer''. What remains open: whether any additional
non-perturbative cascade structure operates dynamically during phase~C
(Big-Bang prelude) and, if so, how it affects the tower-growth identification.

\item \textbf{Higher distinguished layers beyond $d = 217$.} Paper~0's
bilinear-invariant structure singles out the pair $(19, 217)$ as the
distinguished thresholds for $\rho_\Lambda$. No further distinguished $d$ for
this observable. But Bott periodicity continues beyond $d = 217$: a Dirac
layer exists at $d = 29$ (suppressed below the $d_1 = 19$ phase transition;
identified with the neutrino mass scale in the cover sheet) and at $d = 37$,
etc. Whether any of these produces a future phase transition in non-$\Lambda$
observables (e.g., a dark sector at $d = 29$ becoming dynamically active at
some future tick) is open.

\item \textbf{Cascade-order-sensitive observables.} Proposition~\ref{prop:ordering}
establishes a strict cascade-internal ordering for Standard Model feature
addition during inflation. Whether any post-Big-Bang observable retains
memory of this ordering (beyond the standard thermal freeze-out ordering) is
open. Candidates include non-thermal relic abundances or CMB non-Gaussian
correlations between modes seeded at different cascade ticks; no concrete
candidate has been identified.

\end{enumerate}

\section{Tier Classification}
\label{sec:tiers}

Following the tier framework used across the cascade series (most explicitly
in Part~V \S11 and Part~IVb \S11), each result of Part~VI is classified below.
Tiers match the series convention: Tier~1 for theorem-level forced results;
Tier~2 for derived numerical predictions with clean structural arguments; Tier~3
for derivations with caveats or provisional inputs; Tier~5 for proposed
extensions or identifications not yet theorem-level.

\begin{center}
\renewcommand{\arraystretch}{1.3}
\begin{tabular}{lcl}
\toprule
Result & Tier & Justification \\
\midrule
Time = tower growth (Def~\ref{def:tower-time}) & 5 &
  Identification, not theorem \\
Paper~II clarifications (\S\ref{sec:paper2-amendments}) & 3 &
  Structural refinements, no Paper~II theorem modified \\
Phase table (Section~\ref{sec:phases}) & 5 &
  Depends on tick-rate and factor-activation rule \\
$N_e = 70.94$ (Proposition~\ref{prop:efolds}) & 5 &
  Depends on tick-rate (i); landing in window suggestive \\
$N = 217$ Big Bang (Theorem~\ref{thm:reheating}) & 5 &
  Mechanism proposed; thermalisation-to-$w=1/3$ step assumed \\
$g_{\mathrm{eff}}^{\mathrm{cascade}} = 106.75$ (Prop~\ref{prop:g_eff}) & 3 &
  Forced by Part~IVa SM projection; $d=29$ alternative gives $112.00$ \\
$T_{\mathrm{RH}} = 7.40 \times 10^{17}$~GeV (Corollary~\ref{cor:TRH}) & 3 &
  Uses cascade-derived $g_{\mathrm{eff}}$; depends on $\Delta\rho$ at $N{=}217$ \\
Gram deficit as perturbation source (Def~\ref{def:gram-amp}) & 4 &
  Source identified; amplitude formula is a Beta-function theorem \\
$P_\delta(d) \sim 1/(8d^2)$ (Prop~\ref{prop:dspectrum}) & 4 &
  Direct consequence of Paper~0 Supplement Thm~15.2 \\
$P_\delta(k) \propto 1/(\ln k)^2$ (Cor~\ref{cor:kspectrum}) & 5 &
  Functional form derived; comparison to $n_s$ fuzzy at current precision \\
$r$ structurally suppressed (Prop~\ref{prop:tensor-provisional}) & 4 &
  Qualitative argument from no-graviton commitment; magnitude open \\
Tensors as classical induced (Prop~\ref{prop:tensor-induced}) & 3 &
  Follows from Paper~II=III Thm~6.1 + Paper~III Lovelock \\
Particle structural timeline (Prop~\ref{prop:ordering}) & 5 &
  Follows from Paper~IVa layer assignments under tower-growth reading \\
Cover-sheet duality (\S\ref{sec:coversheet}) & 3 &
  Dual descriptions of known cascade picture \\
Falsifier list (\S\ref{sec:falsifiers}) & 4 &
  Structural predictions tied to cascade content \\
\bottomrule
\end{tabular}
\end{center}

No result of Part~VI is Tier~1 or Tier~2. The paper's function is to provide
a speculative but concrete cascade-native cosmic history and a battery of
falsifiers; its content is not robust enough to support precision predictions
at the level of Parts~0--V. If the tower-growth identification matures ---
specifically, if the tick rate $\alpha$ is fixed structurally, the tensor $r$
is computed rigorously, and the perturbation freeze-out mechanism is derived
--- relevant results here would promote to Tier~2 and the paper would be
promoted from speculative extension to regular entry in the series.

\begin{remark}[On the honesty of the tier classification]
Part~VI commits to the discipline of the rest of the series: every result
carries its tier, and Tier~5 is explicitly not sufficient for the cascade to
claim a prediction. The ``inflation e-folds land at $70$'' coincidence is
suggestive but not a prediction; the ``$r$ is structurally suppressed''
argument is qualitative, not quantitative. These are correctly flagged.
Readers should not conflate ``interesting speculation that is internally
self-consistent'' with ``derivation that compels belief.''
\end{remark}

\begin{thebibliography}{9}

\bibitem{paper0} [Author], \textit{The Cascade Series --- Part~0: Scale Variance from Orthogonality}, March 2026.

\bibitem{paper0supp} [Author], \textit{The Cascade Series --- Part~0 Supplement: Inter-Layer Coupling and the Independent-Step Correction}, March 2026.

\bibitem{paper1} [Author], \textit{The Cascade Series --- Part~I: The Cosmological Constant from the Observer's Frame}, March 2026.

\bibitem{paper2} [Author], \textit{The Cascade Series --- Part~II: Quantum Mechanics from the Cascade}, March 2026.

\bibitem{paper3} [Author], \textit{The Cascade Series --- Part~III: General Relativity from the Cascade}, March 2026.

\bibitem{paper23} [Author], \textit{The Cascade Series --- Part II\,=\,III: Quantum Gravity without Quantising Gravity}, March 2026.

\bibitem{paper4a} [Author], \textit{The Cascade Series --- Part IVa: Standard Model from the Cascade}, March 2026.

\bibitem{paper4b} [Author], \textit{The Cascade Series --- Part IVb: Masses, Couplings, and Precision Predictions}, March 2026.

\bibitem{paper5} [Author], \textit{The Cascade Series --- Part V: Cosmology from the Cascade}, March 2026.

\end{thebibliography}

\end{document}
